\begin{trapbox}
\textbf{误区一：椭圆大系数对应长轴。} 先把右边化为 $1$ 再比较分母。
\begin{center}\begin{tikzpicture}[font=\small,>=stealth]
\draw[black!30,rounded corners=2pt] (0,0) rectangle (14.4,8.2);
\node[anchor=west,font=\small\bfseries] at (.5,7.72) {椭圆先归一化：较小平方项系数，对应较大分母和长轴};
\begin{scope}[shift={(3.6,4.1)},scale=.64]
\draw[black!35,->] (-4,0)--(4,0) node[right] {$x$};\draw[black!35,->] (0,-3.5)--(0,3.5) node[above] {$y$};\draw[blue,very thick] (0,0) ellipse (3 and 2);\fill (-2.2360679775,0) circle (2pt);\fill (2.2360679775,0) circle (2pt);\end{scope}
\node at (3.6,1.15) {$4x^2+9y^2=36$};
\begin{scope}[shift={(10.8,4.1)},scale=.64]
\draw[black!35,->] (-4,0)--(4,0) node[right] {$x$};\draw[black!35,->] (0,-3.5)--(0,3.5) node[above] {$y$};\draw[blue,very thick] (0,0) ellipse (2 and 3);\fill (0,-2.2360679775) circle (2pt);\fill (0,2.2360679775) circle (2pt);\end{scope}
\node at (10.8,1.15) {$9x^2+4y^2=36$};

\node[anchor=west] at (.7,.55) {同尺度；除以 $36$ 后，两图长半轴都是 $3$，焦点轴分别为横轴、纵轴};
\end{tikzpicture}\end{center}

$4x^2+9y^2=36$ 的分母为 $9,4$，长轴横向；$9x^2+4y^2=36$ 则纵向。
二者焦点分别为 $(\pm\sqrt5,0)$、$(0,\pm\sqrt5)$。

\medskip
\textbf{误区二：双曲线沿用椭圆减号，或看较大分母定轴。}
双曲线应看正项变量，且 $c^2=a^2+b^2$；实半轴 $a$ 不保证大于虚半轴 $b$。

\medskip
\textbf{误区三：抛物线平方哪个变量，焦点就在哪条轴。}
恰好应看\textbf{一次项}：$y^2=\pm2px$ 的轴是横轴，$x^2=\pm2py$ 的轴是纵轴。
$p/2$ 是顶点到焦点的距离，$p$ 才是焦点到准线的距离；先读 $2p$ 再除以 $2$。

\medskip
\textbf{误区四：只知焦点轴就照抄原点标准式。}
先检查中心或顶点；平移图已经展示两个不同方程。
圆的两分母相等，不能与两个不同焦点的椭圆混为一谈；零分母也不属于本条标准式。
\end{trapbox}
